%%%%%%%%%%%%%%%%%%%%%%%%%%%%%%%%%%%%%%%%%%%%%%%%%%%%%%%%%%%%%%%%%%%%%%%%%%%%
% CAPÍTULO 13: RIESGOS Y MEDIDAS DE CONTROL
%%%%%%%%%%%%%%%%%%%%%%%%%%%%%%%%%%%%%%%%%%%%%%%%%%%%%%%%%%%%%%%%%%%%%%%%%%%%

\chapter{Riesgos y medidas de control}
\label{ch:riesgos}

La Tabla~\ref{tab:riesgos-control} identifica los principales riesgos del proyecto, su impacto previsto y las medidas de mitigación correspondientes.

\begin{longtable}{p{4.2cm} p{2.3cm} p{8.14cm}}
  \caption{Riesgos y medidas de control.}
  \label{tab:riesgos-control}\\
  \hline
  \textbf{Riesgo} & \textbf{Impacto} & \textbf{Mitigación} \\
  \hline
  \endfirsthead
  \multicolumn{3}{c}{\tablename\ \thetable{} -- continuación}\\
  \hline
  \textbf{Riesgo} & \textbf{Impacto} & \textbf{Mitigación} \\
  \hline
  \endhead
  Alcance excesivo del Gemelo Digital & Alto & Limitar el modelo a usuarios, vehículos, viajes, rutas y estados relevantes. \\
  Exceso de complejidad en optimización & Medio/alto & Comenzar con un \textit{matching} heurístico interpretable. \\
  Dificultad para disponer de datos reales & Alto & Usar datos sintéticos y escenarios reproducibles para la evaluación. \\
  Desalineación entre el DT y la aplicación & Medio & Definir contratos de datos y eventos desde el diseño. \\
  Problemas de tiempo & Alto & Priorizar un MVP funcional antes de abordar IA u optimización avanzada. \\
  Resultados poco medibles & Medio & Definir métricas y escenarios antes de implementar la simulación. \\
  Problemas con el motor de mapas/rutas & Medio & Seleccionar una alternativa de código abierto (OSRM) y disponer de un plan B. \\
  Dependencia de herramientas externas & Bajo & Documentar alternativas y mantener el núcleo desacoplado. \\
  Dificultad para validar con usuarios reales & Medio & Aceptar la validación con datos sintéticos y documentar esta limitación. \\
  Desviación del calendario académico & Alto & Realizar revisiones periódicas con el director y priorizar estrictamente el MVP. \\
  \hline
\end{longtable}
